% © 2026 Ing. David Jaroš — CC BY-NC-ND 4.0
%
% This work is licensed under a Creative Commons Attribution-NonCommercial-NoDerivatives
% 4.0 International License (CC BY-NC-ND 4.0).
%
% License History: Earlier drafts (up to v0.3) were released under CC BY 4.0.
% From v0.4 onward, all material is released under CC BY-NC-ND 4.0 to protect
% the integrity of the theoretical work during ongoing academic development.
%
% See LICENSE.md for full license text.

% UBT-AI-PROVENANCE-BEGIN
% schema: ubt-ai-provenance/v1
% tier: B_machine_verified
% ai_assistance: disclosed
% human_review: machine-verification
% editorial_responsibility: Ing. David Jaroš
% policy: ../../AI_PROVENANCE.md
% notice: Machine-verified against named sources or verifiers; individual attestation is not claimed.
% UBT-AI-PROVENANCE-END

\ifdefined\INCLUDEMODE
\else
  \documentclass[12pt,a4paper]{article}
  \usepackage[a4paper, margin=2.5cm]{geometry}
  \usepackage{amsmath, amssymb, amsthm}
  \usepackage{mathtools}
  \usepackage{hyperref}
  \usepackage{xcolor}
  \usepackage{mdframed}
  \usepackage{booktabs}

  \newtheorem{theorem}{Theorem}[section]
  \newtheorem{lemma}[theorem]{Lemma}
  \newtheorem{proposition}[theorem]{Proposition}
  \newtheorem{corollary}[theorem]{Corollary}
  \theoremstyle{definition}
  \newtheorem{definition}[theorem]{Definition}
  \theoremstyle{remark}
  \newtheorem{remark}[theorem]{Remark}

  \newmdenv[backgroundcolor=yellow!20, linecolor=orange, linewidth=1.5pt]{gapbox}
  \newmdenv[backgroundcolor=green!15, linecolor=green!60!black, linewidth=1.5pt]{resultbox}
  \newmdenv[backgroundcolor=red!10, linecolor=red!60, linewidth=1.5pt]{deadendbox}
  \newmdenv[backgroundcolor=blue!8, linecolor=blue!50, linewidth=1.5pt]{insightbox}

  \title{\textbf{Chirality Derivation, Step 3: Gap~C1 Conditional Resolution}\\
         \large $W^\pm$ Vertex is Not $P_\psi$-Invariant Because\\
         Minimal $S[\Theta]$ Contains No $W_R$ Coupling (Selection Rule)}
  \author{David Jaro\v{s}}
  \date{2026-03-06}

  \begin{document}
  \maketitle

  \begin{abstract}
  We partially close Gap~C1: the $W^\pm$ vertex in the UBT action $S[\Theta]$ is
  not $P_\psi$-invariant, implying parity violation.  The argument is structural:
  $S[\Theta]$ contains $SU(2)_L$ acting by left multiplication and $U(1)_Y$
  acting by right multiplication (proved in \texttt{appendix\_E2\_SM\_geometry.tex}).
  There is \textbf{no} $W_R$ (right-handed $SU(2)$) coupling in the
  \emph{minimal} $S[\Theta]$; this is a \textbf{selection rule / model axiom},
  not yet derived from first principles.
  Under this axiom, the $W^\pm$ covariant derivative term is exclusively
  left-handed and $P_\psi(\mathcal{L}_W) \neq \mathcal{L}_W$.
  Parity violation in weak interactions follows from the $\psi$-circle
  orientation \emph{conditional on the no-$W_R$ axiom}.
  See \texttt{open\_problems.md}, items OP-S3 and OP-S4.
  \end{abstract}
\fi

%──────────────────────────────────────────────────────────────────────────────
\section{Setup and Prior Results}
\label{sec:setup}
%──────────────────────────────────────────────────────────────────────────────

From \texttt{step1\_psi\_parity.tex} (Steps~1--5 of this chain):

\begin{center}
\begin{tabular}{lll}
  \toprule
  \textbf{Statement} & \textbf{Status} & \textbf{Reference} \\
  \midrule
  $P_\psi:\psi\to-\psi$ maps $n\to-n$ & Proved & Def.~1, Step~1 \\
  $\partial_\psi$ anti-commutes with $P_\psi$ & Proved & Lem.~2, Step~1 \\
  $P_\psi$ acts as $\gamma^5$ in $\psi$-sector & Proved & Prop.~3, Step~1 \\
  Preferred orientation: $n>0$ for matter & Proved & Lem.~4, Step~1 \\
  $SU(2)_L$ from left action on $\mathcal{B} = \mathbb{C}\otimes_a\mathbb{H}$ & Proved & App.~E2, §6 \\
  $U(1)_Y$ from right action on $\mathcal{B}$ & Proved & App.~E2, §6 \\
  \textbf{Gap C1}: $W^\pm$ vertex $P_\psi$-odd from $S[\Theta]$ & \textbf{Conditional (this step)} & Below \\
  \bottomrule
\end{tabular}
\end{center}

%──────────────────────────────────────────────────────────────────────────────
\section{The Gauge Structure of $S[\Theta]$}
\label{sec:gauge_structure}
%──────────────────────────────────────────────────────────────────────────────

\subsection{Left and Right Actions on $\mathcal{B} = \mathbb{C}\otimes_a\mathbb{H}$}
\label{subsec:LR_actions}

The biquaternion algebra $\mathcal{B} = \mathbb{C}\otimes_a\mathbb{H}
\cong \mathrm{Mat}(2,\mathbb{C})$ admits left and right multiplication
by elements of $\mathcal{B}^*$:
\begin{equation}
  L_q: M \mapsto q\cdot M, \qquad
  R_q: M \mapsto M\cdot q.
\end{equation}
The automorphism group decomposes as:
\begin{equation}
  \mathrm{Aut}(\mathcal{B})
  \;\cong\;
  \bigl[\mathrm{GL}(2,\mathbb{C}) \times \mathrm{GL}(2,\mathbb{C})\bigr]\big/\mathbb{Z}_2,
\end{equation}
giving rise to $\mathrm{Inn}(\mathbb{H})_L \times \mathrm{Inn}(\mathbb{H})_R
\cong SU(2)_L \times SU(2)_R$ (see \texttt{appendix\_E2\_SM\_geometry.tex~§5}).

\subsection{UBT Gauge Content}
\label{subsec:UBT_gauge}

The UBT kinetic action is:
\begin{equation}
  S_{\mathrm{kin}}[\Theta]
  \;=\; \frac{1}{2}\int_{\mathcal{M}\times\mathbb{T}_\psi}
         \bigl\langle D_\mu\Theta,\, D^\mu\Theta \bigr\rangle\,d^4x\,d\psi,
\end{equation}
where the covariant derivative incorporates:
\begin{enumerate}
  \item $SU(2)_L$ gauge field $W^a_\mu$ via \textbf{left} multiplication:
        \begin{equation}
          D_\mu\Theta \;\supset\; -i\,g_2\,W^a_\mu\,T^a\,\Theta,
          \quad T^a = \frac{i\sigma^a}{2}\mathbf{1}_R.
          \label{eq:SU2L_covariant}
        \end{equation}
  \item $U(1)_Y$ gauge field $B_\mu$ via \textbf{right} multiplication:
        \begin{equation}
          D_\mu\Theta \;\supset\; -i\,g_1\,B_\mu\,Y\,\Theta,
          \quad Y = -i\,\mathbf{1}_L.
          \label{eq:U1Y_covariant}
        \end{equation}
\end{enumerate}

\begin{theorem}[Minimal UBT action: $SU(2)_L$ and $U(1)_Y$ only]
\label{thm:no_WR}
The \emph{minimal} UBT action $S[\Theta]$ gauges $SU(2)_L$ (left action)
and $U(1)_Y$ (right scalar phase action).  No $SU(2)_R$ gauge field
$W^a_{R,\mu}$ is introduced.  The absence of $SU(2)_R$ is a
\textbf{selection rule / model axiom} of the minimal theory; it is not
yet derived from first principles.
\end{theorem}

\begin{proof}
The right action of $\mathrm{Inn}(\mathbb{H})_R \cong SU(2)_R$ on
$\mathcal{B}$ would require gauging the map $\Theta \mapsto \Theta\cdot q$
for $q\in SU(2)_R$.  From \texttt{appendix\_E2\_SM\_geometry.tex~§6}
(Theorem~6.1 and Theorem~6.3):
\begin{itemize}
  \item The \emph{left} $SU(2)$ generators are $T^a: M\mapsto (i\sigma^a/2)M$;
        their commutator $[T^a, T^b] = \varepsilon^{abc}T^c$ is proved.
        These are gauged with coupling $g_2$.
  \item The \emph{right} action in the minimal theory is restricted to the
        \emph{scalar phase} $\Theta \mapsto e^{-i\theta}\Theta$,
        giving $U(1)_Y$ with coupling $g_1$.
  \item No $SU(2)_R$ generators are introduced in the covariant derivative.
\end{itemize}
While $\mathcal{B} \cong \mathrm{Mat}(2,\mathbb{C})$ admits a non-scalar
right $SU(2)$ action in principle, this action is not gauged in the
minimal UBT construction.  The reason why gauging $SU(2)_R$ is excluded
--- whether by algebraic obstruction, phenomenological constraint, or
a deeper symmetry-breaking mechanism --- is an open problem.
See \texttt{open\_problems.md}, item OP-S4.
\end{proof}

\begin{gapbox}
\textbf{Open: derivation of $SU(2)_R$ absence.}
The claim that $S[\Theta]$ contains no $SU(2)_R$ coupling is a
\emph{selection rule of the minimal action}, not an algebraic theorem.
A full derivation showing why $SU(2)_R$ cannot be gauged, or why it
decouples, is required before this can be promoted to a proved result.
Until then, the parity-violation argument of this section is
\textbf{conditional} on this axiom.
\end{gapbox}

%──────────────────────────────────────────────────────────────────────────────
\section{Gap C1: $W^\pm$ Vertex is $P_\psi$-Odd}
\label{sec:gap_C1}
%──────────────────────────────────────────────────────────────────────────────

\begin{theorem}[Gap~C1 Conditional Result: $W^\pm$ vertex breaks $P_\psi$]
\label{thm:gap_C1}
Under the no-$W_R$ selection rule (Theorem~\ref{thm:no_WR}), the
$W^\pm$ interaction vertex in $S[\Theta]$ is not $P_\psi$-invariant:
\begin{equation}
  P_\psi \bigl(W^a_\mu\,T^a\,\Theta\bigr)
  \;\notin\; S[\Theta].
\end{equation}
Consequently $P_\psi(\mathcal{L}_W) \neq \mathcal{L}_W$ and parity is
explicitly broken in the $W^\pm$ sector.
\end{theorem}

\begin{proof}
We follow the three-step argument outlined in the problem statement.

\textbf{Step~1.} The $W^\pm$ interaction term is:
\begin{equation}
  \mathcal{L}_{W} \;=\; g_2\,W^a_\mu\,J^{\mu,a}_L,
  \quad
  J^{\mu,a}_L \;=\; \bar{\Theta}_L\,\gamma^\mu\,T^a\,\Theta_L,
\end{equation}
where $\Theta_L$ is the left-handed (odd-winding) sector defined in
Step~1.  The coupling is \emph{exclusively} to $\Theta_L$; there is
no $\Theta_R$ coupling (Theorem~\ref{thm:no_WR}).

\textbf{Step~2.} Under $P_\psi$: winding number $n \mapsto -n$,
so $\Theta_L$ (odd $n$) maps to $\Theta_R$ (odd $n$ with sign change),
i.e., to the right-handed sector.  More precisely (Lemma~2, Step~1):
\begin{equation}
  P_\psi\,\Theta_L \;=\; \Theta_R, \qquad
  P_\psi\,\Theta_R \;=\; \Theta_L.
\end{equation}

\textbf{Step~3.} Under $P_\psi$, the current transforms as:
\begin{equation}
  P_\psi\bigl(\bar{\Theta}_L\,\gamma^\mu\,T^a\,\Theta_L\bigr)
  \;=\; \bar{\Theta}_R\,\gamma^\mu\,T^a\,\Theta_R.
\end{equation}
But $T^a$ acts on the \emph{left-handed} subspace $\mathcal{H}_-$
(the $n > 0$ sector).  Acting on $\mathcal{H}_+ = P_\psi(\mathcal{H}_-)$
(the right-handed sector) gives zero, because $T^a = 0$ on $\mathcal{H}_+$
(by Theorem~\ref{thm:no_WR}: there is no $SU(2)_R$ gauge coupling).
Therefore:
\begin{equation}
  P_\psi\bigl(J^{\mu,a}_L\bigr)
  \;=\; \bar{\Theta}_R\,\gamma^\mu\,T^a\,\Theta_R.
\end{equation}
This right-handed current does \emph{not} appear in $S[\Theta]$ because
no $SU(2)_R$ coupling is introduced in the minimal UBT action
(Theorem~\ref{thm:no_WR}).  Hence:
\begin{equation}
  P_\psi(\mathcal{L}_W) \;\neq\; \mathcal{L}_W.
\end{equation}
Parity is therefore \textbf{explicitly broken}: the $P_\psi$-transformed
action contains a term absent from the original action.

\begin{gapbox}
\textbf{Note on eigenvalue language.}
We do \emph{not} claim $P_\psi(\mathcal{L}_W) = -\mathcal{L}_W$ (a
$(-1)$-eigenvalue statement).  A term that maps to zero under a symmetry
operator is \emph{not} thereby in the $(-1)$-eigenspace: zero lies in
every eigenspace.  The correct statement is simply that
$P_\psi(\mathcal{L}_W) \neq \mathcal{L}_W$, which is sufficient to
establish that $P_\psi$ is not a symmetry of the action.
\end{gapbox}
\end{proof}

\begin{remark}
The key logical chain is:
\begin{center}
\textit{No $W_R$ in $S[\Theta]$ (selection rule)}
$\Longrightarrow$
\textit{$T^a = 0$ on $\mathcal{H}_+$}
$\Longrightarrow$
\textit{$P_\psi(\mathcal{L}_W) = $ right-handed coupling $\notin S[\Theta]$}
$\Longrightarrow$
\textit{$P_\psi(\mathcal{L}_W) \neq \mathcal{L}_W$: parity explicitly broken.}
\end{center}
\end{remark}

%──────────────────────────────────────────────────────────────────────────────
\section{Consequence: Parity Violation from $\psi$-Circle Orientation}
\label{sec:consequence}
%──────────────────────────────────────────────────────────────────────────────

\begin{corollary}[Parity violation from UBT geometry --- conditional]
\label{cor:parity}
Under the no-$W_R$ selection rule, parity violation in weak interactions
(the $V-A$ structure) follows from the $\psi$-circle orientation in UBT.
This is a \textbf{candidate result}, not a theorem, until the no-$W_R$
axiom is derived from first principles.
\end{corollary}

\begin{proof}
By Lemma~4 of Step~1: the preferred orientation $n > 0$ for matter is
fixed by CPT and the arrow of time.  This is an algebraic consequence of
the $\psi$-circle structure.

By Theorem~\ref{thm:gap_C1} (conditional on no-$W_R$): the $W^\pm$ vertex
couples exclusively to odd-winding modes ($n > 0 = $ left-handed), and
$P_\psi(\mathcal{L}_W) \neq \mathcal{L}_W$.  Therefore parity is broken
in the $W^\pm$ sector by the geometry of the $\psi$-circle.

The $V - A$ current $\bar{\psi}\gamma^\mu(1-\gamma^5)\psi/2$ arises
directly: $1 - \gamma^5 \leftrightarrow 1 - P_\psi$ (projection onto
odd-winding modes), which is the UBT left-projector.
\end{proof}

%──────────────────────────────────────────────────────────────────────────────
\section{Status Update}
\label{sec:status}
%──────────────────────────────────────────────────────────────────────────────

\begin{resultbox}
\textbf{Gap~C1: CONDITIONALLY RESOLVED.}

\textbf{Conditional result:} ``The $W^\pm$ vertex is $P_\psi$-odd because
$S[\Theta]$ contains no $W_R$ coupling --- $U(1)_Y$ is the only right-handed
gauge interaction in the \emph{minimal} UBT action.  Therefore $SU(2)_L$
chirality follows from $\psi$-circle orientation, \emph{conditional on the
no-$W_R$ selection rule}.''

\medskip
\textbf{Complete chirality derivation chain (with status):}
\begin{enumerate}
  \item $\mathbb{H}$ gives $SU(2)_L\times SU(2)_R$ from left/right actions
        --- Proved (App.~E2, §5).
  \item $P_\psi$ acts as $\gamma^5$ in $\psi$-sector --- Proved (Step~1).
  \item Preferred orientation $n>0$ by CPT --- Proved (Step~1, Lem.~4).
  \item $T^a = 0$ on $\mathcal{H}_+$ (no $W_R$ in minimal action)
        --- \textbf{Model axiom / selection rule} (Thm.~\ref{thm:no_WR}).
        Derivation open (OP-S4).
  \item $W^\pm$ vertex causes $P_\psi(\mathcal{L}_W)\neq\mathcal{L}_W$
        --- \textbf{Candidate} (conditional on item~4).
  \item Parity violation from $\psi$-circle --- \textbf{Candidate}
        (Cor.~\ref{cor:parity}; conditional on items~4--5).
\end{enumerate}

\medskip
\textbf{Conditions on which the above depends:}
\begin{itemize}
  \item No $SU(2)_R$ gauging in the minimal action (selection rule, not yet derived).
  \item Grade/psi-mode equivalence (Gap C1 in \texttt{open\_problems.md}, OP-S2).
  \item Antiunitary structure of $T$ (OP-S1).
\end{itemize}

\medskip
\textbf{Consequence (conditional):} The weak interaction is parity-violating
because the $\psi$-circle has a preferred orientation and $SU(2)$ acts only
from the left in the minimal action.

\textbf{Update DERIVATION\_INDEX.md:} ``$SU(2)_L$ chirality (not $SU(2)_L\times
SU(2)_R$)'' $\to$ \textbf{Candidate [L0]} (conditional on no-$W_R$ axiom)
--- see this file.
\end{resultbox}

\ifdefined\INCLUDEMODE
\else
  \end{document}
\fi
